\documentclass[a4paper,11pt,twoside,onecolumn,openright,final,oldfontcommands]{memoir} % final / draft

\usepackage[portuges]{babel}% portuguese  portuges
%\addto\captionsportuges{\renewcommand{\contentsname}{\'{I}ndice}}% alternative to %\renewcommand{\contentsname}{\'{I}ndice} inside document

\usepackage[utf8]{inputenc}
\usepackage[T1]{fontenc}

\usepackage{amsmath,amsfonts,amstext,amssymb,amsthm} % Utilidades matematicas
\usepackage[]{graphics} % option: draft

\usepackage{subfig}% subfigure
\usepackage{caption}
%\usepackage{keyval}

\usepackage{rotating}

\usepackage{stmaryrd}% brackets in jump function

%\usepackage{psfrag} % [scanall] option to include \tex commands in the eps files
\usepackage{color}
%\usepackage{bbm}
%\usepackage[Symbol]{upgreek}%Euler Symbol Symbolsmallscale

\usepackage{kbordermatrix}

% \usepackage{tikz}
% \usetikzlibrary{shapes.misc,positioning,arrows}


\usepackage{arydshln}% array dash line
\setlength{\dashlinedash}{3pt} %for the length of each dash segment in a dash line
\setlength{\dashlinegap}{3pt} %controls the amount of each gap between dash segments


%\usepackage{bm,upgreek}
%\makeatletter
%\newcommand{\bfgreek}[1]{\bm{\@nameuse{up#1}}}
%\makeatother

% norma portuguesa ... indent first word
\usepackage{indentfirst}

% Translineacao de palavras com particulas
% para funcionar inserir "- no ponto, como em foi"-se
% \defineshorthand{"-}{\nobreak\hskip0pt\discretionary{-}{-}{-}\nobreak\hskip0pt}



%\usepackage{calrsfs} % change \mathcal fonts
%\usepackage{calligra}

% \usepackage{mathrsfs}  use as $\mathscr{ABCDEF}$
%
 \makeatletter
\providecommand*{\DisemulatePackage}[1]{%
  \@namelet{ver@#1.\@pkgextension}\relax}
\makeatother \DisemulatePackage{index}

\usepackage{multirow}

\renewcommand{\citeindexfile}{authors}
\makeindex[\citeindexfile]
\usepackage{natbib}
\citeindextrue
%
\makeindex
\bibpunct{(}{)}{,}{a}{,}{,}



\usepackage{longtable}
%\usepackage{dpfloat}
%\usepackage{ifthen}
%\usepackage{afterpage}

%\DeclareFontFamily{T1}{calligra}{}
%\DeclareFontShape{T1}{calligra}{m}{n}{<->s*[1.3]callig15}{}

\usepackage{fancybox}


%\DeclareRobustCommand\renewcommand\ss{\fontfamily{calligra}\fontencoding{T1}\selectfont}
%\newcommand\textcalligra[1]{\begingroup\renewcommand\ss{\symbol{255}\kern-.22em}\fontfamily{calligra}\fontencoding{T1}\selectfont#1\endgroup}

\DeclareMathAlphabet{\mathbbmsl}{U}{bbm}{m}{sl}
\DeclareMathAlphabet{\mathbbmbxsl}{U}{bbm}{bx}{sl}

%vertical
\setlength{\voffset}{-0.5cm} \setlength{\topmargin}{0cm}
\setlength{\headheight}{6mm} \setlength{\headsep}{1cm}
\setlength{\textheight}{23cm} \setlength{\footskip}{0.8cm}

%horizontal
\setlength{\hoffset}{-1in} \setlength{\marginparwidth}{0cm}
\setlength{\evensidemargin}{2.5cm} \setlength{\textwidth}{15cm}
\setlength{\marginparsep}{0cm} \setlength{\oddsidemargin}{3.0cm}



% quote marks
%\newcommand{\quotationmarks}[1]{``#1''}

% upright div
%\newcommand{\udiv}{\mathrm{div} \,}
%\newcommand{\udiag}{\mathrm{diag} }
%\newcommand{\usym}{\text{sym}}

% upright d
%\newcommand{\ud}{\,\mathrm{d}}
% \renewcommand{\bs}[1]{\boldsymbol{#1}}
%\newcommand{\bs}[1]{\boldsymbol{#1}}

%\newcommand{\abs}[1]{\lvert#1\rvert}
%\newcommand{\norm}[1]{\lVert#1\rVert}
%\newcommand{\rank}[1]{\text{rank}\left(#1\right)}
%\newcommand{\jump}[1]{\left\llbracket #1 \right \rrbracket}
%\newcommand{\degree}{\ensuremath{^\circ}}


% Definitions of products
%\newcommand{\tensorialp}[2]{#1\otimes#2}
%\newcommand{\vectorialp}[2]{#1\times#2}
%\newcommand{\dotp}[2]{#1\cdot#2}
%\newcommand{\scalarp}[2]{#1 : #2}

%\newcommand{\sen}{\text{sen}}

\setlength{\defaultaddspace}{5pt} % distance for table spacing

%\theoremstyle{remark}
%\newtheorem{rem}{Remark}[section]

\newenvironment{mydescription}%
{\begin{list}%
{}%
{\renewcommand{\makelabel}[1]{\textbf{\emph{##1}}}}}%
{\end{list}}


%
% Numerate until subsections
\maxsecnumdepth{subsection}
%
% write until subsections in ToC
\maxtocdepth{subsection}
%

%\definecolor{mygrey}{gray}{0.6}
\pagestyle{companion}

%
% Here starts the document
%
\begin{document}
%
%\renewcommand{\contentsname}{\\'{I}Índice} % it must go here, after \begin{document} command
%
\frontmatter % for starting numbering i, ii, iii, iv,... and no numbering of sectional units
%

%
\thispagestyle{empty}%
%

\begin{minipage}[c]{0.2\textwidth}
    \centering
\hspace{-2cm}\includegraphics[width=2.5cm,clip]{figures/frontmatter/Logo_IST_color.eps}
\end{minipage}%
\begin{minipage}[c]{0.8\textwidth}
    \centering
\begin{Large}
\textbf{
UNIVERSIDADE TÉCNICA DE LISBOA\\[0.5cm]
INSTITUTO SUPERIOR TÉCNICO}
\end{Large}
\end{minipage}\\[3cm]%
%
\mbox{}
\vfill
\begin{center}
\textbf{\huge{Análise não-linear de estruturas utilizando elementos finitos tridimensionais}}\\[3cm]
\textbf{\Large{Rui Maciel}}\\
\vfill
\today
\vfill
\end{center}
%
\cleardoublepage % does not write in the back
%\movetooddpage % does write in the back
%

%
%\movetooddpage
\setcounter{page}{1}
%
\tableofcontents
%
\movetooddpage
%
\listoffigures
%
\movetooddpage
%
\listoftables
%
\movetooddpage
%
%\include{symbols}
%
%\movetooddpage
%
%\include{abbreviations}
%
%\movetooddpage
%
\mainmatter

%
\chapter{Introdução}

\cleardoublepage

%
\chapter{O problema de elasticidade tridimensional}

\cleardoublepage
                                                                                                                                                                                               
%
\chapter{O método dos elementos finítos}

\cleardoublepage
                                                                                                                                                                                                    
%
\chapter{Discretização em elementos-fronteira}

\cleardoublepage
                                                                                                                                                                                            
%
\chapter{Implementação}

\cleardoublepage
                                                                                                                                                                                            
%
\chapter{Exemplos}

\cleardoublepage
                                                                                                                                                                                                
%
\chapter{Conclusão}

\cleardoublepage
 

%\movetooddpage
%
% References
%
%GATHER{bibdata.bib}
\bibliography{bibdata}
\bibliographystyle{bibliotiago}
%
% Citation index
%
%\renewcommand{\indexname}{Citations index}
%\printindex[\citeindexfile]
%
% Subject index
%
%\renewcommand{\indexname}{Subject index}
%\printindex
%
\end{document}
